% =============================================================================
% TERMO DE CONSENTIMENTO LIVRE E ESCLARECIDO (TCLE)
% Versão para Especialistas (Painel Delphi / Comitê de Adaptação)
% =============================================================================

\documentclass[12pt,a4paper,brazil]{article}

\usepackage[utf8]{inputenc}
\usepackage[T1]{fontenc}
\usepackage[brazil]{babel}
\usepackage{geometry}
\usepackage{setspace}
\usepackage{indentfirst}
\usepackage{graphicx}
\usepackage{fancyhdr}

\geometry{a4paper, left=3cm, right=2cm, top=3cm, bottom=2cm}
\onehalfspacing

\pagestyle{fancy}
\fancyhf{}
\rhead{\small TCLE --- Especialistas}
\lhead{\small PPGPI/UFS}
\cfoot{\thepage}

\begin{document}

\begin{center}
{\large\bfseries UNIVERSIDADE FEDERAL DE SERGIPE}\\[0.3cm]
{\normalsize Programa de Pós-Graduação em Ciência da Propriedade Intelectual (PPGPI)}\\[1cm]
{\Large\bfseries TERMO DE CONSENTIMENTO LIVRE E ESCLARECIDO (TCLE)}\\[0.3cm]
{\normalsize Versão para Especialistas (Painel Delphi e Comitê de Adaptação Transcultural)}\\[0.5cm]
\noindent\rule{\textwidth}{0.4pt}
\end{center}

\vspace{0.5cm}

\noindent\textbf{Título da Pesquisa.} Modelo Computacional para Salvaguarda de Saberes Agrícolas Tradicionais como Ativos Intangíveis em Comunidades Quilombolas do Semiárido Nordeste~II.

\noindent\textbf{Pesquisadora Responsável.} Catuxe Varjão de Santana Oliveira.

\noindent\textbf{Orientador.} Prof. Dr. Paulo Roberto Gagliardi (PPGPI/UFS).

\noindent\textbf{Coorientador.} Prof. Dr. Luiz Diego Vidal Santos (UEFS).

\noindent\textbf{Instituição.} Universidade Federal de Sergipe (UFS), Programa de Pós-Graduação em Ciência da Propriedade Intelectual.

\vspace{0.5cm}

\section*{Convite à Participação}

V.Sa. está sendo convidado(a) a participar, na condição de especialista, de uma pesquisa de doutorado cujo objetivo é desenvolver e validar uma ferramenta computacional que auxilie comunidades quilombolas na salvaguarda dos seus próprios saberes agrícolas tradicionais. Este documento contém as informações necessárias para sua decisão informada sobre a participação. Solicita-se a leitura integral antes de manifestar seu consentimento.

\section*{Objetivo da Pesquisa}

A pesquisa tem como objetivo geral desenvolver e validar um modelo computacional integrado, articulando adaptação transcultural de instrumentos, elicitação estruturada de variáveis, modelagem psicométrica e lógica fuzzy, para que comunidades quilombolas do Território de Identidade Semiárido Nordeste~II (Bahia) disponham de uma ferramenta auditável de avaliação e priorização da salvaguarda de seus saberes e sistemas agrícolas tradicionais. A ferramenta resultante ficará à disposição das comunidades para uso em processos junto a IPHAN, INPI, CGen e outros órgãos. Sua participação como especialista contribuirá para a validação de conteúdo dos instrumentos e para a convergência de variáveis relevantes para a avaliação de vulnerabilidade biocultural.

\section*{Descrição da Participação}

Dependendo do componente em que V.Sa. for convidado(a) a participar, as atividades previstas serão as seguintes.

\paragraph{Comitê de Adaptação Transcultural (Componente~A, Etapa~4).} Avaliação dos itens do questionário WOCAT-SLM adaptado em quatro dimensões de equivalência (semântica, idiomática, experiencial e conceitual), utilizando escala de 4 pontos. O formulário de avaliação será enviado por meio eletrônico, com prazo de 15 dias para preenchimento. O tempo estimado de participação é de 60 a 90 minutos.

\paragraph{Painel Delphi (Componente~B).} Participação em três rodadas iterativas de consulta estruturada, com intervalos de 30 a 45 dias entre rodadas. A Rodada~1 envolve enumeração aberta de variáveis relevantes para avaliação de vulnerabilidade e salvaguarda de saberes agrícolas tradicionais. A Rodada~2 consiste em avaliação de variáveis consolidadas em escala Likert de 5 pontos (relevância, clareza e operacionalidade). A Rodada~3 compreende reavaliação com feedback agregado anonimizado. O tempo estimado por rodada é de 30 a 45 minutos, totalizando aproximadamente 2 horas ao longo de 4 meses.

\paragraph{Grupo Focal Comparativo (Componente~B).} Participação em sessão de grupo focal com 8 a 10 especialistas, com duração estimada de 90 a 120 minutos, gravada em áudio. A sessão abordará as mesmas variáveis do Delphi em formato de discussão coletiva, para fins de comparação metodológica entre técnicas estruturadas e não estruturadas.

\section*{Riscos}

Os riscos da participação são mínimos. Embora não se esperem desconfortos significativos, a divergência de opiniões durante o processo Delphi ou grupo focal pode gerar algum desconforto pontual. V.Sa. poderá interromper sua participação a qualquer momento, sem necessidade de justificativa. Os dados serão apresentados de forma anonimizada, sem atribuição individual de posicionamentos. Para painelistas Delphi que não completarem ao menos 2 das 3 rodadas, os dados parciais serão excluídos da análise, sem qualquer ônus para o(a) participante.

\section*{Benefícios}

A participação contribuirá para o avanço metodológico na integração entre adaptação transcultural, elicitação estruturada e modelagem computacional aplicada a contextos de comunidades tradicionais. O instrumento resultante (WOCAT-SLM-QBR) poderá ser utilizado por pesquisadores, gestores e instituições para documentação de conhecimentos tradicionais em outros contextos. Os mestres de saberes quilombolas que integrarem o comitê receberão reconhecimento como coautores do instrumento adaptado. Todos os participantes receberão cópia dos resultados consolidados e do instrumento final validado.

\section*{Sigilo e Proteção de Dados}

As respostas individuais serão codificadas e tratadas de forma anonimizada. No processo Delphi, o feedback entre rodadas será apresentado de forma agregada, sem identificação de posicionamentos individuais. Os dados serão armazenados em ambiente digital seguro e criptografado, com acesso restrito à equipe de pesquisa, e mantidos por 5 anos após publicação dos resultados, conforme Resolução CNS nº~466/2012. Os áudios de grupo focal serão destruídos após transcrição e verificação. O tratamento de dados pessoais seguirá a Lei Geral de Proteção de Dados (Lei nº~13.709/2018).

\section*{Participação Voluntária e Direito de Desistência}

A participação é voluntária. V.Sa. pode recusar-se a participar ou desistir a qualquer momento, sem justificativa e sem qualquer prejuízo profissional ou acadêmico. A recusa ou desistência não gerará ônus de qualquer natureza.

\section*{Contatos}

Em caso de dúvidas sobre a pesquisa, entre em contato com a pesquisadora responsável, Catuxe Varjão de Santana Oliveira, pelo telefone (\rule{1.5cm}{0.4pt}) \rule{3cm}{0.4pt} ou pelo e-mail \rule{5cm}{0.4pt}.

Em caso de denúncia ou dúvida sobre questões éticas, entre em contato com o Comitê de Ética em Pesquisa da UFS (CEP/UFS), localizado na Cidade Universitária Prof. José Aloísio de Campos, Av. Marechal Rondon, s/n, Jardim Rosa Elze, São Cristóvão/SE, CEP 49100-000, pelo telefone (79) 3194-7208 ou pelo e-mail cepufs@academico.ufs.br.

\vspace{1cm}

\section*{Consentimento do(a) Participante}

Declaro que li este Termo de Consentimento Livre e Esclarecido, que compreendi as informações nele contidas e que concordo voluntariamente em participar desta pesquisa. Recebi uma via deste documento e tive oportunidade de fazer perguntas, que foram respondidas de forma satisfatória.

\vspace{1.5cm}

\noindent
\begin{minipage}[t]{0.55\textwidth}
Local e data:\\[1.5cm]
\rule{8cm}{0.4pt}
\end{minipage}

\vspace{1cm}

\noindent
\begin{minipage}[t]{0.48\textwidth}
\centering
\rule{7cm}{0.4pt}\\
Nome e assinatura do(a) participante
\end{minipage}
\hfill
\begin{minipage}[t]{0.48\textwidth}
\centering
\rule{7cm}{0.4pt}\\
Nome e assinatura do(a) pesquisador(a)
\end{minipage}

\vspace{1cm}

\noindent
\footnotesize{Este documento é emitido em duas vias de igual teor, ficando uma com o(a) participante e outra com a pesquisadora responsável.}

\end{document}
